\subsection{True intrinsic convexity on the same original prescribed-class line}

\begin{definition}[The genuine section on the original polarized model line]
\label{local:actual-original-polarized-model-generic-section}
\lean{EnvelopeProperty.ProjSubvariety.normalizedSimultaneousOriginalRatioSectionAt}
\leanok
\uses{local:actual-normalized-polynomial-original-ratio-line-iso}
At every actual normalized-model point with nonzero original first ideal value,
pull the constructed whole actual polarization comparison to its original
residue field and apply it to the genuine polynomial homogeneous section
tensored with the actual canonical ratio section. This gives an actual generic
section of the SAME accepted original line $p^*\mathcal O_X(d)\otimes L_{I/J}$.
\end{definition}
\begin{proof}
The preceding whole line comparison is an actual isomorphism on the same
original normalized model. Its actual residue-field pullback is therefore an
actual map between the corresponding generic line modules. Apply its true
top-section linear map to the already constructed genuine tensor section.
No original section or metric comparison is supplied as additional data.
\end{proof}

\begin{definition}[The true weight on the original prescribed-class line]
\label{local:actual-original-polarized-model-section-weight}
\lean{EnvelopeProperty.ProjSubvariety.normalizedSimultaneousOriginalRatioWeightAt}
\leanok
\uses{local:actual-original-polarized-model-generic-section}
The actual original polarized weight is the negative logarithm of the genuine
model norm of that actual generic section in the SAME accepted original line.
Its model norm uses the true valuative center on the same actual proper model.
\end{definition}
\begin{proof}
The actual original line is invertible by its already proved tensor-factor
invertibility. The actual normalized-model structure is proper over the
original polynomial base over any field, and the original Gauss valuation
is bounded on every original coefficient. The genuine center-defined model
norm is therefore available for this actual section; apply its negative logarithm.
\end{proof}

\begin{lemma}[The constructed whole comparison is a genuine all-point model isometry]
\label{local:actual-original-polarized-model-weight-isometry}
\lean{EnvelopeProperty.ProjSubvariety.normalizedSimultaneousOriginalRatioWeightAt_eq}
\leanok
\uses{local:actual-original-polarized-model-section-weight,local:actual-normalized-polynomial-original-ratio-line-iso}
At EVERY actual normalized-model point, the preceding TRUE original-line weight
equals the TRUE polynomial-polarized section weight. This equality includes the
original residue valuation and all zero-section logarithmic values.
This isometry itself is valid over any original field; perfectness is not
needed for the whole-line isometry or its negative-logarithm equality.
\end{lemma}
\begin{proof}
The genuine whole-line-isomorphism model norm theorem is valid for every
original invertible-line isomorphism: pulling it through the actual valuative
center preserves the original center frames, hence their entire original
generic section norms. Apply it to the constructed whole polarization
comparison and the genuine original tensor section. Its image is exactly the
section defined above. Apply the negative logarithm to this actual norm equality.
\end{proof}

\begin{theorem}[Original local convexity gives true intrinsic convexity on the original line]
\label{local:actual-prescribed-class-original-model-intrinsic-convexity}
\lean{EnvelopeProperty.ProjSubvariety.IsLocallyConvexAll.exists_prescribedClass_intrinsicConvex_originalModelMetric}
\leanok
\uses{local:actual-prescribed-class-intrinsic-convex-normal-model,local:actual-original-polarized-model-weight-isometry}
For the original projectively normal variety over any field, positive
rational class $a$, and original locally $\mathcal O(a)$-convex real function
$\varphi$, there exist positive $r,d$ with EXACT original $d=ra$, an integer
$C$, and actual original ideals with everywhere finite original product Gauss
logarithms, on their same actual normal proper model, such that the true actual
ratio weight is $r\varphi-C$ at every actual model point. For every original
admissible degree-$d$ form, the TRUE weight on the accepted ORIGINAL line
$p^*\mathcal O_X(d)\otimes L_{I/J}$ is intrinsically convex PL in genuine
original regular units on the entire actual model form-open. No line-comparison,
metric, outgoing-slope, numerical-degree, or nefness certificate is assumed.
\end{theorem}
\begin{proof}
The established all-point intrinsic model-convexity theorem constructs the
same original $r,d,C,I,J$, retaining $d=ra$, the whole actual ratio-weight
formula, and genuine regular-unit convexity of every actual polynomial-base
polarized section weight. Retain these constructed objects and their exact
original ratio identity. For each original admissible form, the preceding
isometry identifies its TRUE polynomial-line weight with the TRUE original-line
weight on every point of the same actual form-open. The genuine convexity
congruence theorem transfers the actual regular-unit presentations along that
pointwise equality. This yields the claim on the SAME accepted original line,
without changing the original degree, ratio model, point locus or coefficient
valuation.
\end{proof}
